%%%%%%%%%%  scan idx 700  |  folio 687  %%%%%%%%%%
\SGAPrintPage{687}%
that is isolated are geometric invariants that are by no means trivial, and will undoubtedly have their part to play in the ``arithmetic'' study of these local rings.

Of course, once one has a reasonable notion of a Chow ring for a regular noetherian scheme admitting an ample invertible Module, one should generalize to this case the relations established or conjectured, respectively, in n\textsuperscript{o} 4 between this ring and the ring $\Gtop^\bullet(X)$, and formulate in terms of this ring the question of the ``Riemann--Roch formula without denominators'' stated in n\textsuperscript{o} 3.

\sgaStableTarget{sga6.e14.bibliography-of-expose-xiv}\subsection*{BIBLIOGRAPHY}
\begin{enumerate}[label={[\arabic*]}]
\item M. F. Atiyah and F. Hirzebruch, The Riemann--Roch theorem for analytic embeddings, Topology, vol. 1, p. 151-166 (1962).
\item A. Borel and J.P. Serre, Le théorème de Riemann--Roch, Bull. Soc. Math. France, vol. 86, p. 97-136 (1958).
\item R. Bott and H. Samelson, Amer. Journ. of Math. vol 80, p. 1004 (1958).
\item C. Chevalley, Les classes d'équivalence rationnelle, I and II, Seminar at the ENS 1958, exposés 2 and 3.
\item C. Chevalley, Classification des groupes de Lie algébriques, Seminar at the ENS 1956/58.
\item A. Dold and D. Puppe, Homologie nicht-additiver Funktoren. Anwendungen, Annales de l'Inst. Fourier t. XI, p. 201-312 (1961).
\item H. Grauert, Ein Theorem der analytischen Garbentheorie, Pub. Math. n\textsuperscript{o} 5, p. 5-64 (1960).
\item A. Grothendieck, Classes de faisceaux et théorème de Riemann--Roch, mimeographed notes prepared by the Institute for Advanced Study, 1957 (work cited as [RRR] in the present exposé, and reproduced in our seminar as an Appendix to Chap. O).
\item A. Grothendieck, La théorie des classes de Chern, Bull. Soc. Math. France vol. 86, p. 137-154 (1958).
%%%%%%%%%%  scan idx 701  |  folio 688  %%%%%%%%%%
\SGAPrintPage{688}%
\item A.
Grothendieck, Classes de Chern et représentations linéaires des groupes discrets, in: Dix exposés sur la cohomologie des schémas, Masson--North Holland Pub. Cie
\item H. Hironaka, Resolution of singularities of an algebraic variety over a field of characteristic zero, Ann. of Math. 79, p. 109-326 (1964).
\item L. Hodgkin, thesis (to be published?).
\item L. Illusie, work in preparation.
\item J.P. Serre, Algèbre Locale, Multiplicités, Lecture Notes in Math. n\textsuperscript{o} 11 (1965) Springer.
\item A. Borel, Sous-groupes commutatifs et torsion des groupes de Lie compacts connexes, Tohoku Math. Journal, vol. 13, p. 216-240 (1961).
\item[{[16]}(*)] P. Deligne, Exp. XVIII in Théorie des topos et cohomologie étale des schémas (SGA 4), by M. Artin, A. Grothendieck and J. L. Verdier.
\item[{[17]}(*)] U. Manin, Correspondances, motifs et transformations monoïdales (in Russian). Mat. Sbornik, t 77 (119) n\textsuperscript{o} 4, 1968.
\end{enumerate}

\SGAAuthorFootText{*}{Added in February 1969}

%%%%%%%%%%  scan idx 702  |  folio 689  %%%%%%%%%%
\SGAPrintPage{689}%
\sgaStableTarget{sga6.e14.note-for-page-1}\subsection*{Note for page 1.}

The theory of the operations $\Lambda^i$ in $D^-(X)$, and of the operations $\lambda^i$ in $K^\bullet(X)$, was written down by P. Deligne in 1968. He proves that $K^\bullet(X)$ thereby becomes a pre-$\lambda$-ring, and, when $\underline{O}_X$ is a sheaf of Algebras over a field, that $K^\bullet(X)$ is even a $\lambda$-ring. It is plausible that a refinement of Deligne's method would make it possible to dispense with the restrictive hypothesis on $\underline{O}_X$.

\sgaStableTarget{sga6.e14.notes-for-page-11}\subsection*{Notes for page 11}

\noindent(*) (April 1969) D. MUMFORD points out to us a proof of (4.6) in the \uline{Chow ring} (without disregarding torsion), which he obtained in a seminar at Harvard in 1959! The reader will find the proof of this beautiful result in SGA 5 VII 9. This result obviously contains the corresponding cohomological result mentioned in note (**). However, the cohomological result is proved by J. P. JOUANOLOU under more general conditions (without a quasi-projectivity hypothesis, and in the relative smooth case over an arbitrary base).

\medskip

\noindent(**) (January 1969) J. P. JOUANOLOU has just proved these formulas for every $y$ without disregarding torsion, cf. SGA 5 VII \S{} 4.9. DELIGNE and MANIN 17 had previously observed independently, and by (easy) arguments of a rather different nature, that it was enough to establish formulas for ``algebraic'' $y$, and even for $y=1$ in 16, to conclude the same relation in the general case.

% Final-index helper macros
\providecommand{\Qcoh}{\operatorname{Qcoh}}
\providecommand{\Coh}{\operatorname{Coh}}
\providecommand{\Parf}{\operatorname{Parf}}
\providecommand{\Lib}{\operatorname{Lib}}
\providecommand{\Loclib}{\operatorname{Loclib}}
\providecommand{\Proj}{\operatorname{Proj}}
\providecommand{\Mod}{\operatorname{Mod}}
\renewcommand{\Ob}{\operatorname{Ob}}
\providecommand{\Hom}{\operatorname{Hom}}
\providecommand{\toramp}{\operatorname{tor.amp}}
\providecommand{\parfamp}{\operatorname{par.amp}}
\providecommand{\coh}{\operatorname{coh}}
\providecommand{\qcoh}{\operatorname{qcoh}}
\providecommand{\parf}{\operatorname{parf}}
\providecommand{\torf}{\operatorname{torf}}
\providecommand{\amp}{\operatorname{amp}}
\providecommand{\num}{\operatorname{num}}
\providecommand{\naif}{\operatorname{naif}}
\providecommand{\alg}{\operatorname{alg}}
\providecommand{\Pic}{\operatorname{Pic}}
\providecommand{\Gr}{\operatorname{Gr}}
\providecommand{\Fil}{\operatorname{Fil}}
\providecommand{\Ch}{\operatorname{Ch}}
\providecommand{\Todd}{\operatorname{Todd}}
\providecommand{\Diff}{\operatorname{Diff}}
\providecommand{\Dcat}{\mathrm D}
\providecommand{\Tr}{\operatorname{Tr}}
\providecommand{\Z}{\mathbb Z}
\providecommand{\Q}{\mathbb Q}

\newpage
%%%%%%%%%%  scan idx 703  |  folio 690  %%%%%%%%%%
\SGAPrintPage{690}%

